\documentclass[10pt]{article}
\providecommand{\FIGDIR}{fig}
\newcommand{\DIGESTDATE}{18 September 2026}
% ===================== PM Research Watch - shared preamble =====================
\usepackage[T1]{fontenc}
\usepackage[utf8]{inputenc}
\usepackage{amsmath,amssymb}
\usepackage{tgpagella}
\usepackage{tgheros}
\renewcommand{\sfdefault}{qhv}
\renewcommand{\ttdefault}{lmtt}
\usepackage{microtype}
\usepackage[english]{babel}

\usepackage[a4paper,top=17mm,bottom=17mm,left=17mm,right=17mm,
            headsep=6mm,footskip=10mm]{geometry}
\usepackage{graphicx}
\usepackage{xcolor}
\usepackage{colortbl}
\usepackage{tikz}
\usetikzlibrary{calc}
\usepackage[most]{tcolorbox}
\usepackage{titlesec}
\usepackage{fancyhdr}
\usepackage{booktabs}
\usepackage{array}
\usepackage{longtable}
\usepackage{enumitem}
\usepackage{multicol}
\usepackage{ragged2e}
\usepackage[hidelinks]{hyperref}

% ------------------------------------------------------------------ palette
\definecolor{Ink}    {HTML}{14262E}
\definecolor{Deep}   {HTML}{0F4C5C}
\definecolor{Teal}   {HTML}{1F7A8C}
\definecolor{Sky}    {HTML}{4FA3C4}
\definecolor{Sage}   {HTML}{6FA287}
\definecolor{Amber}  {HTML}{E3A008}
\definecolor{Clay}   {HTML}{B5643C}
\definecolor{Coral}  {HTML}{D1495B}
\definecolor{Violet} {HTML}{6A5B9E}
\definecolor{Slate}  {HTML}{4A6572}
\definecolor{Paper}  {HTML}{FBF9F5}
\definecolor{Mist}   {HTML}{EEEAE2}
\definecolor{Rule}   {HTML}{D9D3C9}

\pagecolor{Paper}
\color{Ink}

\hypersetup{colorlinks=true, linkcolor=Deep, urlcolor=Teal, citecolor=Teal,
            pdftitle={PM Research Watch}, pdfauthor={Automated daily digest}}

% ------------------------------------------------------------------ headings
\titleformat{\section}
  {\sffamily\bfseries\large\color{Deep}}{}{0pt}
  {}[{\vspace{1.2mm}\color{Amber}\titlerule[1.1pt]}]
\titlespacing*{\section}{0pt}{5.5mm}{2.8mm}

\titleformat{\subsection}
  {\sffamily\bfseries\normalsize\color{Teal}}{}{0pt}{}
\titlespacing*{\subsection}{0pt}{3.6mm}{1.4mm}

\setlist[itemize]{leftmargin=4.2mm, itemsep=1.1mm, topsep=1.3mm, parsep=0pt,
                  label=\textcolor{Teal}{\footnotesize$\blacktriangleright$}}

% ------------------------------------------------------------------ header/footer
\pagestyle{fancy}
\fancyhf{}
\renewcommand{\headrulewidth}{0pt}
\renewcommand{\footrulewidth}{0pt}
\fancyhead[L]{\sffamily\scriptsize\color{Slate}PM RESEARCH WATCH\;\textcolor{Rule}{\textbar}\;\DIGESTDATE}
\fancyhead[R]{\sffamily\scriptsize\color{Slate}Particulate matter monitoring \& health effects}
\fancyfoot[C]{\sffamily\scriptsize\color{Slate}\thepage}

% ------------------------------------------------------------------ components
% metric tile
\newcommand{\metric}[3]{%
  \begin{tcolorbox}[enhanced, colback=white, colframe=Rule, boxrule=0.5pt,
    arc=1.6mm, left=2.2mm, right=2.2mm, top=1.8mm, bottom=1.8mm,
    borderline west={1.7mm}{0pt}{#3}, nobeforeafter, valign=center,
    height=17mm, width=\linewidth]
    \RaggedRight
    {\sffamily\bfseries\fontsize{17}{17}\selectfont\color{#3}#1}\\[0.5mm]
    {\sffamily\fontsize{6.5}{7.4}\selectfont\color{Slate}#2}
  \end{tcolorbox}}

% pull-quote / key insight box
\newtcolorbox{keybox}[1][]{enhanced, colback=Mist, colframe=Mist,
  arc=1.6mm, left=3.2mm, right=3.2mm, top=2.4mm, bottom=2.4mm,
  borderline west={1.7mm}{0pt}{Amber}, #1}

% subtopic band
\newcommand{\band}[2]{%
  \vspace{3.2mm}\noindent
  \begin{tcolorbox}[enhanced, colback=#2, colframe=#2, arc=1.2mm,
    left=2.6mm, right=2.6mm, top=1.1mm, bottom=1.1mm, boxrule=0pt,
    width=\linewidth, nobeforeafter]
    {\sffamily\bfseries\small\color{white}#1}
  \end{tcolorbox}\vspace{1.6mm}}

% one paper entry: {short cite}{title}{journal}{doi}{body}
\newcommand{\paperentry}[5]{%
  \par\vspace{1.9mm}\noindent
  {\sffamily\bfseries\footnotesize\color{Ink}#1}\;%
  {\sffamily\fontsize{7}{8}\selectfont\color{Slate}#3}\\[0.4mm]
  {\footnotesize\itshape\color{Deep}#2}\\[0.6mm]
  {\footnotesize #4}\\[0.3mm]
  {\fontsize{6.6}{7.6}\selectfont\color{Slate}\texttt{doi:}\,\href{https://doi.org/#5}{\texttt{\detokenize{#5}}}}
  \par}

% caption style
\newcommand{\figcap}[1]{\par\vspace{1.0mm}
  {\fontsize{7.2}{8.6}\selectfont\color{Slate}\textsf{\textbf{Fig.}}\; #1\par}\vspace{1.5mm}}

\providecommand{\permil}{\textperthousand}

\setlength{\parindent}{0pt}
\setlength{\parskip}{0pt}
\renewcommand{\arraystretch}{1.18}


\begin{document}

% ==================================================================== MASTHEAD
\thispagestyle{empty}
\begin{tikzpicture}[remember picture, overlay]
  \fill[Deep] (current page.north west) rectangle
      ($(current page.north east)+(0,-67mm)$);
  \fill[Amber] ($(current page.north west)+(0,-67mm)$) rectangle
      ($(current page.north east)+(0,-68.4mm)$);
  \foreach \x/\y/\r in {22/12/0.9, 41/26/1.5, 63/9/0.7, 88/31/1.1, 112/17/1.9,
                        138/29/0.8, 159/13/1.3, 178/54/1.0, 196/21/1.6,
                        30/44/0.6, 74/57/1.2, 125/48/0.9, 168/6/0.8, 191/59/0.7,
                        52/18/1.0, 100/60/0.6, 148/20/1.1, 66/49/1.4, 205/45/0.8}
    \fill[Sky, opacity=0.30] ($(current page.north west)+(\x mm,-\y mm)$) circle (\r mm);
\end{tikzpicture}

\vspace*{4mm}
{\sffamily\fontsize{7.6}{9}\selectfont\color{Sky}%
 AUTOMATED DAILY LITERATURE DIGEST\;\textbullet\;PhD RESEARCH WATCH\par}
\vspace{2.2mm}
{\sffamily\bfseries\fontsize{27}{29}\selectfont\color{white}%
 Particulate Matter\par}
{\sffamily\fontsize{27}{29}\selectfont\color{Sky}%
 Monitoring \& Health Effects\par}
\vspace{5.0mm}
{\sffamily\fontsize{8.4}{10}\selectfont\color{white}%
 Issue 2026-09-18\;\;\textcolor{Amber}{\textbar}\;\;Entry window 18 September 2026 (full day)\;\;%
 \textcolor{Amber}{\textbar}\;\;30 records screened in scope\par}
\vspace{18mm}

% -------------------------------------------------------------- metric strip
\noindent
\begin{minipage}[t]{0.190\linewidth}\metric{30}{RECORDS IN SCOPE\\(102 EXCLUDED)}{Deep}\end{minipage}\hfill
\begin{minipage}[t]{0.190\linewidth}\metric{8}{OF 10 SUBTOPIC\\BANDS POPULATED}{Teal}\end{minipage}\hfill
\begin{minipage}[t]{0.190\linewidth}\metric{7}{EFFECT ESTIMATES\\(ONLY 3 WITH A CI)}{Sky}\end{minipage}\hfill
\begin{minipage}[t]{0.190\linewidth}\metric{5}{TIER-A\\RECORDS}{Amber}\end{minipage}\hfill
\begin{minipage}[t]{0.190\linewidth}\metric{8}{RECORDS CARRIED ON\\METADATA ALONE}{Clay}\end{minipage}

\vspace{6mm}


% -------------------------------------------------------------- signal of day
\section*{\sffamily\bfseries\color{Deep}Signal of the day}
\vspace{-1.5mm}{\color{Amber}\rule{\linewidth}{1.1pt}}\vspace{2.5mm}

\begin{keybox}
{\sffamily\bfseries\small\color{Deep}Four independent records today attack the same
assumption from four directions: that a surface-sited, mass-reporting monitor measures the
exposure that matters. Source, height, composition and speciated dose each break it.}\\[1.4mm]
\footnotesize
\textbf{Source.} Zhang et al.\ (\textit{Nat Commun}) separated wildfire-specific from
non-wildfire PM$_{2.5}$ across \textbf{20 US states, 2006--19}, in a self-controlled
design. Per 1\,$\mu$g/m$^3$, wildfire PM$_{2.5}$ raised cardiopulmonary hospitalisation
risk \textbf{10\% (heart failure) to 16\% (asthma)}; non-wildfire PM$_{2.5}$ raised it
\textbf{4.7\% (COPD) to 8.5\% (hypertension)}. Equal mass, unequal hazard, by roughly a
factor of two.
\textbf{Height.} Bai et al.\ (\textit{Atmos Meas Tech}) used Doppler wind lidar over
Hefei to show two dust episodes with opposite surface signatures: in one, a descending
low-level jet drove aerosol down and surface PM$_{10}$ peaked at
\textbf{410\,$\mu$g/m$^3$}; in the other, a stable jet core held the layer aloft and the
surface peak arrived \textbf{$\sim$8\,h} after frontal onset. A surface network reads
clean air while a loaded layer sits overhead.
\textbf{Composition.} Yu et al.\ (\textit{Atmos Chem Phys}) found marine organic-aerosol
light absorption over the Yellow and Bohai Seas rising with both organic-nitrogen
concentration \emph{and} ON/OC ratio --- absorption capability per unit carbon changes
with continental outflow, which is precisely the term a fixed optical calibration assumes
constant.
\textbf{Speciated dose.} Morrel et al.\ (\textit{Biol Psychiatry Glob Open Sci}), the
largest human study in the issue at \textbf{6,949 ABCD preadolescents}, finds the signal
in \emph{components} and not in mass: zinc, ammonium and bromine predicted perivascular-%
space count; calcium, zinc and ammonium predicted poorer cognition (FDR $p<0.01$); and
frontal-lobe perivascular-space count \textbf{mediated} the zinc--cognition association.
Total PM$_{2.5}$ entered only as a weighted-quantile-sum mixture, and never as the
quantity a monitor reports.
Only one of the four is a sensing paper. Together they are the strongest argument this
watch has carried that the next increment in exposure assessment is source-, height- and
composition-resolved, not denser mass.
\end{keybox}

\vspace{3mm}


% -------------------------------------------------------------- exec brief
\section*{\sffamily\bfseries\color{Deep}What changed today}
\vspace{-1.5mm}{\color{Amber}\rule{\linewidth}{1.1pt}}\vspace{2.5mm}

\footnotesize
\begin{itemize}
\item \textbf{A three-year reference-grade absorption record at a rural South Asian site.}
Raju et al.\ (\textit{J Environ Sci Health A}) report seven-wavelength AE33 aethalometry at
Panchgaon, Haryana, \textbf{January 2020 -- December 2022}: eBC \textbf{4.06\,$\pm$\,2.15},
fossil-fuel \textbf{3.10\,$\pm$\,2.40} and biomass-burning \textbf{1.42\,$\pm$\,1.19}\,%
$\mu$g/m$^3$, with eBCff carrying \textbf{68\%} of absorption at 880\,nm and winter maxima
roughly double the monsoon minima. This is the co-location anchor low-cost networks in the
region do not have. The apportionment itself is shakier than the record: the
eBCff/eBCbb ratio is \textbf{3.02\,$\pm$\,3.80} --- a standard deviation larger than the
mean --- and the assumed absorption \AA ngstr\"om exponents are never stated.
\item \textbf{A community-designed monitoring campaign that is a template, not a result.}
Frueh et al.\ (\textit{J Expo Sci Environ Epidemiol}) ran \textbf{38 VOCs at 20 sites for
12 months} in a South Philadelphia fenceline community with Philly Thrive, around a former
refinery. Benzene --- the pollutant the community named --- came in low
(\textbf{median 1.21\,$\mu$g/m$^3$}), while (\textit{m,p})-xylenes emerged in industrial
areas, and PMF apportioned sources. One-week integrated passive samplers resolve none of
the short plumes residents actually report, which is the design's own limit. The siting
logic --- community-set locations, a declared pollutant of concern, apportionment rather
than compliance --- is what distributed PM networks should copy.
\item \textbf{A humidity threshold that lands squarely on the sensor-correction problem.}
Sun et al.\ (\textit{Environ Sci Technol}) split wintertime basin nitrate formation almost
evenly between daytime OH+NO$_2$ (\textbf{52\%}) and nighttime N$_2$O$_5$ hydrolysis
(\textbf{46\%}), and locate abrupt efficiency thresholds at \textbf{RH 87\%} and
\textbf{ALWC 76.5\,$\mu$g/m$^3$}. That is the same humidity band in which optical low-cost
sensors diverge from reference monitors --- the aerosol is physically changing where the
instrument is least trustworthy. One site, one winter, no out-of-sample validation.
\item \textbf{A cardiovascular case-crossover in which every PM term is null.}
Chiang et al.\ (\textit{J Am Heart Assoc}) found no same-day association for any pollutant
across \textbf{798 acute aortic dissections}; only CO at 2-day lag survived
(\textbf{OR 2.13 per IQR, 95\% CI 1.075--4.239}). PM$_{2.5}$, PM$_{10}$, NO$_2$, O$_3$ and
SO$_2$ showed nothing at any lag. With six pollutants across four lags and no multiplicity
correction, the surviving association is fragile --- but the PM nulls are themselves
informative: exposure was assigned from \textbf{24 fixed stations}, a density at which
within-city contrast is largely erased.
\item \textbf{Eight of thirty records assert no finding at all.} Metadata-only is the
largest single design bucket in this issue, ahead of measurement campaigns at seven.
Same-day deposits from Elsevier, ACS, BMC and now Taylor \& Francis carry a title, authors
and a DOI and nothing else. A new recovery route was tested and \textbf{failed}: the
OpenAlex per-work endpoint (\texttt{/works/doi:}), which is free even though the
date-range filter is not, returns \texttt{abstract\_inverted\_index} only where the
publisher deposited an abstract to Crossref in the first place --- so it recovers nothing
these publishers withhold.
\end{itemize}

\begin{keybox}
{\sffamily\bfseries\footnotesize\color{Deep}Run note --- this issue was rebuilt}\\[1.2mm]
{\fontsize{7.2}{8.8}\selectfont
A first build of this issue ran at \textbf{10:33 CDT} and shipped 17 records against
roughly ten hours of a twenty-three-hour entry window, contrary to the 22:00 rule the
pipeline adopted on 19 August. Re-harvesting the same window at \textbf{23:0x} returned
\textbf{PubMed 12} (was 3), \textbf{Europe PMC 46} (was 44) and \textbf{Crossref-by-ISSN 53}
(was 24); 13 of the 72 fresh candidates were admitted. This issue supersedes the morning
build and covers the whole of 18 September. The earlier version was never the complete
day, and no record it carried has been removed.\par}
\end{keybox}

\vspace{3mm}


% ==================================================================== ANALYTICS
\section*{\sffamily\bfseries\color{Deep}Corpus analytics}
\vspace{-1.5mm}{\color{Amber}\rule{\linewidth}{1.1pt}}\vspace{3mm}

\noindent\includegraphics[width=\linewidth]{\FIGDIR/f1_subtopics.png}
\figcap{Subtopic assignment of the 30 in-scope records; eight of the ten bands are
populated. \textit{Exposure assessment \& modelling} is the largest at nine, and with
\textit{Sensing, forecasting \& instrumentation} at five the measurement side holds
\textbf{14 of 30}. \textit{Mechanistic toxicology} also reaches five, then
\textit{Cardiovascular \& metabolic} and \textit{Occupational \& indoor} at three each,
\textit{Respiratory \& allergic} and \textit{Neuro / mental health} at two, and
\textit{Burden, policy \& mitigation} at one. \textit{Reproductive \& developmental} and
\textit{Other clinical endpoints} are empty.}

\vspace{2mm}
\noindent\includegraphics[width=\linewidth]{\FIGDIR/f2_design_endpoint.png}
\figcap{Left --- study architecture. \textbf{Metadata only is the largest single bucket at
eight of thirty}, ahead of measurement campaigns at seven; that ordering, not the total,
is the finding, and it is a property of same-day publisher deposits rather than of the
literature. Primary observational epidemiology contributes seven (four acute, two
cross-sectional, one cohort), evidence synthesis four, modelling two and experimental
toxicology two. Right --- health endpoint, canonicalised. \textbf{``No health endpoint''
is 16 of 30} --- the aerosol-chemistry, emissions, source-policy and not-applicable
records all collapse into it --- so \textbf{more than half the issue reports no human
outcome}. Respiratory leads the clinical endpoints at four; cardiovascular, neuro /
cognitive, other clinical and toxicological (in vitro) carry two each.}

\vspace{2mm}
\noindent\includegraphics[width=\linewidth]{\FIGDIR/f3_metric_geography.png}
\figcap{Left --- particle metric. \textbf{Unspeciated PM or emissions terms cover 14 of
30}, and the composition of that bar is the point: it absorbs both the metadata-only
records and the four VOC / precursor studies, none of which measured a particle mass at
all. PM$_{2.5}$ alone accounts for eight, speciated composition and joint PM$_{2.5}$ +
PM$_{10}$ three each, PM$_{10}$ alone and optical properties one each. Right ---
geographic provenance. The \textit{Global / multi-region} bar is ten and again is not what
it looks like: \textbf{four are records whose study location is not recoverable}, and of
the remaining six, three are reviews or commentaries with no field site. China is the
largest genuine single-country bar at six, then the USA at four.}

\vspace{2mm}
\noindent\includegraphics[width=\linewidth]{\FIGDIR/f6_lifecourse.png}
\figcap{Life-course windows addressed; counts are non-exclusive. Working-age adulthood
carries five --- the UK Biobank COPD cohort, the 18-study cement-dust worker
meta-analysis, the Ghanaian mining communities, the Budapest sleep panel and the 798
aortic-dissection cases. Puberty and adolescence is newly occupied this issue by the
\textbf{6,949 ABCD preadolescents}, and older adults by the 130 Seoul patients with
idiopathic pulmonary fibrosis. The paediatric entry is the Toluca series of 266 children
under five; because bronchiolitis cases there are necessarily infants but no infant
stratum is tabulated, the prenatal/infant window is recorded as zero rather than inferred.}

\vspace{2mm}
\noindent\includegraphics[width=\linewidth]{\FIGDIR/f4_heatmap.png}
\figcap{Subtopic $\times$ study-architecture cross-tabulation. The structural fact is
unchanged by tripling the sensing row and doubling the modelling row: \textbf{the
\textit{Exposure assessment \& modelling} row (nine records) and the \textit{Sensing}
row (five) canonicalise to ``no health endpoint'' without a single exception}.
Fourteen records --- the entire measurement side of the issue --- report no human
outcome, and no record anywhere in the matrix links an instrument to a clinical endpoint.
The second feature is dispersion: metadata-only entries appear in six of the eight
populated rows, so the abstract-free deposit problem is not confined to one literature.}

\vspace{2mm}
\noindent\includegraphics[width=\linewidth]{\FIGDIR/f5_forest.png}
\figcap{All seven extractable effect estimates on a common logarithmic axis; only three
report a confidence interval. Those three now bracket the null in both directions ---
two protective hazard ratios from the UK Biobank greenness analysis, and one
\textbf{OR 2.13 (95\% CI 1.075--4.239)} for acute aortic dissection per IQR of CO at
2-day lag --- while \textbf{all four estimates above the null that come from PM exposure
are point values with no interval reported}. Exposure contrasts are not mutually
comparable: hazard ratios per greenness category at a 1000\,m buffer, risk ratios per
1\,$\mu$g/m$^3$ of source-resolved PM$_{2.5}$, an odds ratio per 0.19\,ppm of CO. Read
direction and precision, never relative magnitude.}

\clearpage

% ==================================================================== DIGEST
\section*{\sffamily\bfseries\color{Deep}Paper-level digest}
\vspace{-1.5mm}{\color{Amber}\rule{\linewidth}{1.1pt}}\vspace{1mm}

{\fontsize{7.4}{9}\selectfont\color{Slate}Ordered by cluster. Each entry gives the
methodological core and the finding that matters, not an abstract paraphrase. Records
marked \textit{metadata only} carry no recoverable abstract and assert no finding.\par}

\band{Neuro, cognition \& mental health\hfill 2 records}{Deep}

\paperentry{Temporao et al.\ 2026}
{Daily associations between outdoor air pollution and sleep quality}
{medRxiv (preprint)}
{Budapest Sleep, Traits and Experiences Study --- seven days per participant of daily
activity and sleep-quality reports plus home EEG giving sleep efficiency, \%REM, onset
latency and total sleep time. Linear mixed models with participant random intercepts,
adjusted for age, sex, BMI, education, weekday, season, temperature, physical activity,
screen exposure and travel. \textbf{Day-1 ozone predicted lower day-2 sleep efficiency
and poorer perceived quality}, surviving multiple-comparison correction, multi-pollutant
models and sensitivity analyses for depression, insomnia and chronotype;
\textbf{no robust association for PM$_{2.5}$ or NO$_2$}. The null is confounded with
exposure-model quality --- city-aggregated station data give near-zero within-city
contrast, and misclassification is worse for PM$_{2.5}$ than for the more spatially
homogeneous ozone. Sample size is not stated in the abstract.}
{10.64898/2026.09.15.26363130}

\paperentry{Morrel et al.\ 2026}
{Outdoor air pollution, perivascular space morphology, and cognition in preadolescence}
{Biol Psychiatry Glob Open Sci}
{Cross-sectional analysis of \textbf{6,949 ABCD participants aged 9--10}, with annual-average
PM$_{2.5}$, O$_3$, NO$_2$ and \textbf{15 PM$_{2.5}$ components} assigned to residential
addresses by spatiotemporal models, perivascular-space (PVS) count and volume from T1- and
T2-weighted MRI, and cognition from the NIH Toolbox. \textbf{Zinc, ammonium and bromine}
were positively associated with regional PVS count; higher count in five regions tracked
poorer cognition; \textbf{calcium, Zn and NH$_4^+$} were associated with poorer cognition
(FDR-corrected $p<0.01$); frontal-lobe PVS count \textbf{mediated} the Zn--total-cognition
association ($p<0.01$); and weighted-quantile-sum PM$_{2.5}$ mixtures predicted higher
temporal and cingulate PVS with poorer working memory and crystallised intelligence. The
design cannot order exposure, PVS and cognition, and the abstract reports significance
without a single effect size. The component-specific pattern is the measurement claim:
speciation, not mass, carries the signal.}
{10.1016/j.bpsgos.2026.100790}

\band{Cardiovascular \& metabolic\hfill 3 records}{Teal}

\paperentry{Zhang et al.\ 2026}
{Cardiopulmonary hospitalisation risks from wildfire-specific and non-wildfire
PM$_{2.5}$ in 20 US states}
{Nat Commun}
{Self-controlled design over 20 states, 2006--19, contrasting two-year average
wildfire-specific against non-wildfire PM$_{2.5}$. Per 1\,$\mu$g/m$^3$, wildfire
PM$_{2.5}$ carried \textbf{10\% (heart failure) to 16\% (asthma)} higher hospitalisation
risk; non-wildfire PM$_{2.5}$ \textbf{4.7\% (COPD) to 8.5\% (hypertension)}. Associations
were stronger among minorities, metropolitan residents, the less educated and more
deprived communities. The toxicity contrast at equal mass is the contribution, and it
undercuts the mass-only metric every regulatory network and low-cost sensor reports. Two
limits: no confidence interval is given for any of the eight headline risks, so precision
cannot be judged; and wildfire-specific PM$_{2.5}$ is itself modelled, so the toxicity
difference is conditional on a smoke-attribution step whose error is not propagated.}
{10.1038/s41467-026-76974-7}

\paperentry{Chiang et al.\ 2026}
{Ambient carbon monoxide as a delayed trigger of acute aortic dissection: a time-stratified
case-crossover study in a tropical industrial metropolis}
{J Am Heart Assoc}
{Time-stratified case-crossover of \textbf{798 emergency-department-confirmed acute aortic
dissections}, 2018--2021, exposures assigned from \textbf{24 monitoring stations},
conditional logistic regression across lags 0--3 adjusted for meteorology. \textbf{No
pollutant was associated on the event day.} At 2-day lag each 1\,ppm CO increment carried a
2.67-fold risk (95\% CI 1.207--5.900) and each IQR increment of 0.19\,ppm raised risk
\textbf{113\% (95\% CI 1.075--4.239)}; \textbf{PM$_{2.5}$, PM$_{10}$, NO$_2$, O$_3$ and
SO$_2$ showed nothing at any lag}. Two problems: a 1\,ppm contrast is roughly five IQRs, so
the headline number extrapolates past the observed range; and with six pollutants across
four lags the one surviving association is not multiplicity-corrected. The PM nulls are
the relevant part here --- station-assigned exposure at this network density returns
nulls that a denser deployment might resolve.}
{10.1161/jaha.126.049936}

\paperentry{Fleming \& Fleming 2026 \textnormal{---} metadata only}
{Tire wear pollution and cardiovascular risk}
{Am J Med}
{Elsevier deposited title and authorship on 18 September with no abstract in Crossref,
Europe PMC core, Europe PMC \texttt{fullTextXML} or the OpenAlex per-work record. The title
places it as a commentary linking tire-wear particle pollution to cardiovascular risk; it
asserts no finding here. Logged because the particle class matters for measurement:
\textbf{non-exhaust emissions now dominate traffic PM in fleets with modern
aftertreatment}, and tire wear differs from combustion PM in both size distribution and
composition --- a combination that mass-only optical sensors and exhaust-derived emission
factors each mis-handle. Carried forward for pickup when an abstract appears.}
{10.1016/j.amjmed.2026.09.003}

\band{Respiratory \& allergic\hfill 2 records}{Sage}

\paperentry{San Juan-Reyes et al.\ 2026}
{Short-term ambient air pollution and respiratory hospitalisation in children under five}
{Environ Sci Pollut Res Int}
{266 children under five admitted to one regional IMSS hospital in the Toluca Valley
Metropolitan Zone during 2023; matched case-crossover with conditional logistic
regression, primary window lag 1--7, adjusted for temperature and humidity. PM$_{2.5}$
was associated with admission for viral pneumonia, bacterial pneumonia and asthma, CO
with bronchiolitis, with heterogeneous municipality-level associations. \textbf{The
abstract reports direction only --- not one odds ratio, interval or unit appears} --- so
the effect size is unverifiable from the published summary. 266 events split across four
diagnoses, five pollutants, several lag windows and three municipalities is thin, and
exposure is entirely fixed-station with no personal or indoor component in a setting
where indoor sources plausibly dominate a young child's dose.}
{10.1007/s11356-026-38237-2}

\paperentry{Chen et al.\ 2026}
{Residential green and blue space, genetic susceptibility and incident COPD}
{Int J Chron Obstruct Pulmon Dis}
{385{,}502 UK Biobank participants, Cox models, greenness at 300\,m and 1000\,m buffers
against a COPD polygenic risk score. At 1000\,m: \textbf{green space HR 0.87
(95\% CI 0.83--0.92)}, \textbf{natural environment HR 0.83 (0.78--0.87)}; blue space weak
and inconsistent. No multiplicative or additive interaction with genetic risk, and the
natural-environment association held in every genetic stratum --- so the protective signal
is not running through a susceptible subgroup. The inferential problem is what is absent:
\textbf{no air-pollution term is reported}. Residential greenness is inversely correlated
with traffic PM$_{2.5}$ and NO$_2$ nearly everywhere, so an unadjusted greenness--COPD
association may simply restate a pollution--COPD association.}
{10.2147/copd.s627458}



\band{Mechanistic toxicology\hfill 5 records}{Violet}

\paperentry{Gao et al.\ 2026}
{CCL22-associated immune signatures in air pollution-related atopic dermatitis}
{Front Public Health}
{Shared pollutant--atopic dermatitis targets enriched in cytokine and chemokine
signalling, pattern-recognition receptor activity, IL-17 and Toll-like receptor pathways.
A framework over 15 algorithms and 175 combinations selected plsRglm\,+\,AdaBoost, with
AUC 0.963 in training and 1.000, 0.909 and 0.966 across three validation cohorts, giving
a ten-gene signature (CCL22, CCL5, CSF2, F2RL1, HRH4, ICAM1, IFNG, IL10, IL17A, IL18).
CCL22 was upregulated in lesional samples, localised to dendritic cells and macrophages,
and confirmed in DNFB-induced murine dermatitis. \textbf{The exposure link is the weak
joint: no particulate matter was measured or administered anywhere}, so air pollution
enters only as a database annotation. An AUC of exactly 1.000 indicates leakage or
overfitting rather than performance, and DNFB is a contact hapten, not a particle.}
{10.3389/fpubh.2026.1906298}

\paperentry{Vera-Col\'on et al.\ 2026 \textnormal{---} metadata only}
{Airborne fine particulate matter and human embryonic stem cell osteogenic
differentiation}
{Part Fibre Toxicol}
{BMC deposited this DOI on 18 September with no abstract; Crossref, Europe PMC core,
Europe PMC \texttt{fullTextXML} and the publisher landing page all returned nothing on
per-DOI retry --- \textbf{the first time this gap has appeared on a BMC rather than an
Elsevier deposit}. Carried on title because skeletal development is an endpoint the PM
literature has barely touched, and a human-embryonic-stem-cell osteogenic model would be
a new mechanistic line rather than another airway or endothelial assay. No finding is
asserted: neither particle source, dose, duration nor the direction of the reported
disruption is verifiable from metadata. On the retry list.}
{10.1186/s12989-026-00701-9}

\paperentry{Dulong et al.\ 2026}
{Effects of fine particulate matter on neuroblastoma cell line: metabolic and lipid
deregulations and their consequences on the onset of neurological diseases}
{Biochimie}
{Narrative review built around the \textbf{SH-SY5Y} human neuroblastoma line as an in vitro
neurodegeneration model. The proposed cascade runs mitochondrial damage
$\rightarrow$ ER stress $\rightarrow$ oxidative stress $\rightarrow$ calcium dyshomeostasis
$\rightarrow$ lipid dysregulation, converging on apoptosis and autophagy-dependent
\textbf{ferroptosis}, with lipid peroxidation advanced as the emerging driver. No pooled
estimate, no exposure--response and no dose metric: the review does not harmonise
PM$_{2.5}$ mass, source or exposure duration across the studies it assembles, and the in
vitro doses it draws on generally sit far above ambient, so the chain is qualitative. It
concedes its own in-vitro-to-human translation gap. Every mechanism it privileges is
\textbf{composition-dependent} --- redox-active metals, organic carbon --- which a
mass-only network cannot resolve.}
{10.1016/j.biochi.2026.09.011}

\paperentry{Wu et al.\ 2026}
{$\beta$-carotene alleviates PM$_{2.5}$-induced pulmonary fibrosis via suppression of the
Frizzled5-mediated Wnt5a/Ca$^{2+}$ pathway}
{Ecotoxicol Environ Saf}
{C57BL/6 mice given \textbf{intratracheal PM$_{2.5}$}, paired with BEAS-2B human bronchial
epithelial cells carrying Fzd5 overexpression. Instillation produced fibrosis, oxidative
stress and calcium dyshomeostasis --- more mitochondria-associated ER membranes, altered
IP$_3$R, STIM1, ORAI1 and SERCA2 --- with activation of a \textbf{Wnt5a/Fzd5/Ca$^{2+}$}
axis. Dietary $\beta$-carotene reversed these changes, and Fzd5 overexpression abolished
the protection, identifying Fzd5 as the mediating receptor. \textbf{No instilled dose,
particle source, treatment concentration, effect size or CI} is reported for any endpoint,
so potency cannot be placed against ambient exposure. Intratracheal instillation also
bypasses the deposition physics that a sensor-derived exposure metric is meant to
approximate.}
{10.1016/j.ecoenv.2026.120816}

\paperentry{Kou et al.\ 2026}
{Pollutants regulate changes in pathological markers of neurodegenerative diseases: a new
perspective in environmental toxicology}
{Front Toxicol}
{Review organised around core neurodegenerative pathological markers ---
\textbf{amyloid-$\beta$, tau, $\alpha$-synuclein, TDP-43} --- rather than around individual
pollutants, detailing how air pollutants, metals and persistent organic pollutants modulate
the generation, aggregation, propagation and toxicity of each, under a proposed
``pollutants--pathogenic links--neural network failure'' framework. No quantitative
synthesis, no effect estimates, no exposure metric and no study-quality appraisal; it
reorganises existing mechanistic claims rather than adding evidence, and \textbf{air
pollutants are treated as a single undifferentiated class} alongside metals and POPs. The
one useful implication for measurement is that marker-level convergence makes pollutant
class, not PM mass, the exposure variable worth resolving.}
{10.3389/ftox.2026.1814795}

\band{Occupational \& indoor exposure\hfill 3 records}{Clay}

\paperentry{Khan et al.\ 2026}
{Cement industry emissions and respiratory health: global systematic review and
meta-analysis, 2000--2025}
{Front Public Health}
{Eighteen observational studies, PROSPERO CRD420251268913. Random-effects pooling gave
standardised mean differences of \textbf{$-$0.71 for FEV$_1$, $-$0.56 for FVC and $-$0.57
for FEV$_1$/FVC} in exposed groups; chronic cough was significantly more common while
other symptoms were elevated but not significant. Occupational exposure produced greater
FEV$_1$ impairment than residential --- the expected dose ordering and the most reassuring
internal check available. Reporting is the weakness: \textbf{no confidence interval, no
$I^2$ and no prediction interval accompanies any pooled estimate}, so neither precision
nor heterogeneity can be assessed, and SMDs pooled across studies using different
spirometry reference equations with no common dust metric cannot be read as a
dose--response.}
{10.3389/fpubh.2026.1887654}

\paperentry{Assan et al.\ 2026}
{Ambient particulate matter, urinary arsenic and resident health in two Ghanaian mining
communities}
{Health Sci Rep}
{65 residents each from an artisanal and an incorporated mining community, all resident
at least five years; PM by Modular Air Quality System, urinary arsenic by AAS, kidney
function by albumin--creatinine ratio. More artisanal residents fell in the at-risk band
for kidney disease ($p = 0.006$); the incorporated community showed poorer general health
status. Descriptive only: $n = 130$, chi-square on dichotomised uACR with no confounder
adjustment and no exposure--response, and \textbf{neither PM concentrations nor arsenic
levels are numerically reported} despite being the two measured exposures. The 35 and
70\,$\mu$g/m$^3$ figures quoted as WHO standards are the 2005-era interim targets, not
the 2021 guideline values of 15 and 45\,$\mu$g/m$^3$ for 24-hour PM$_{2.5}$ and
PM$_{10}$ --- so the exceedance framing understates the gap.}
{10.1002/hsr2.73246}

\paperentry{Ikeda et al.\ 2026 \textnormal{---} metadata only}
{Characterizing layout-dependent breathing-zone contaminant concentration in office spaces
with mixing ventilation using an occupant-to-exhaust distance metric}
{Build Environ}
{Same-day Elsevier deposit; nothing in Crossref, Europe PMC core, Europe PMC
\texttt{fullTextXML} or the OpenAlex per-work record. The title indicates a modelling study
deriving breathing-zone contaminant concentration in mixing-ventilated offices from an
\textbf{occupant-to-exhaust distance metric}. In scope for sensor siting specifically: it
is a claim that indoor concentration is a predictable function of geometry, which is the
assumption every single-point indoor monitor makes implicitly and almost never tests.
Carried forward for pickup.}
{10.1016/j.buildenv.2026.115282}

\band{Sensing, forecasting \& instrumentation\hfill 5 records}{Amber}

\paperentry{Bai et al.\ 2026}
{Low-level jet evolution and vertical aerosol redistribution: a Doppler wind lidar study}
{Atmos Meas Tech}
{Two dust episodes over Hefei, East China. April 2021: the northwesterly jet core
descended from 2.5\,km to 0.3\,km then re-ascended to 1.0\,km while surface PM$_{10}$
peaked at \textbf{410\,$\mu$g/m$^3$}, with shear above \textbf{0.04\,s$^{-1}$} and
increased spectral width at the lower jet interface and a downward shift of enhanced
backscatter. March 2022: a stable jet core and persistent lower-tropospheric stability
held aerosol aloft, the PM$_{10}$ peak arriving at 16:00\,LT, \textbf{$\sim$8\,h after
inferred frontal onset}. Two regimes --- rapid downward mixing during jet descent versus
delayed surface impact from a decoupled layer. Two cases is mechanism illustration, not
climatology, and the inertial-oscillation attribution rests on hodograph rotation alone.
The clearest statement yet of why a surface-only network can read clean for hours.}
{10.5194/amt-19-5951-2026}

\paperentry{Huang et al.\ 2026}
{Interpretable spatiotemporal coupling network for emission prediction in sludge
gasification}
{Environ Sci Technol}
{PISTC-Net couples a directed process topology with per-feature optimal time lags,
reaching $R^2$ of \textbf{91.18\% for SO$_2$, 80.37\% for NO$_x$ and 78.09\% for PM} on
industrial-scale data. Under simulated abrupt sensor failure it continued tracking
emission trajectories via learned process inertia, and hierarchical attention shifted from
upstream thermochemical zones for SO$_2$ to whole-process attention for PM, flagging
disturbance propagation \textbf{5--8 minutes} before emission peaks. That PM is the
worst-predicted species is the informative result: particle formation runs through
nucleation, condensation and agglomeration paths a topology-plus-lag representation
captures least well, whereas gas-phase sulfur is nearly a mass balance. Single plant, no
external validation site, and interpretability is inferred from attention maps, which are
a saliency artefact rather than a causal readout.}
{10.1021/acs.est.6c09135}

\paperentry{Kalluri et al.\ 2026 \textnormal{---} metadata only}
{Long-term (2008--2024) black carbon over a semi-arid rain-shadow region in southern
peninsular India}
{Atmos Environ}
{The expected same-day Elsevier deposit gap: no abstract in Crossref, Europe PMC or
PubMed on per-DOI retry. Carried on title because a 17-year single-site black-carbon
record in a semi-arid rain-shadow setting is the kind of long baseline trend detection
needs and South Asian networks mostly lack, and because absorption-coefficient records
are what constrain the optical-to-mass conversion every filter- and scattering-based
instrument depends on. No trend, magnitude or source attribution is asserted --- the title
names methods, not results. On the retry list.}
{10.1016/j.atmosenv.2026.122366}

\paperentry{Krajnc Poga\v{c}nik et al.\ 2026 \textnormal{---} metadata only}
{Selective real-time monitoring of respiratory droplets for indoor air safety in
educational environments}
{Build Environ}
{Elsevier same-day deposit, abstract absent from Crossref, Europe PMC and PubMed on
per-DOI retry. Retained because selective real-time discrimination of respiratory
droplets against a background of ambient particles is the hard problem for any optical
counter deployed indoors: a scattering sensor reports total number in a size bin and
cannot by itself separate exhaled aerosol from resuspended dust, so any claimed
selectivity mechanism is worth reading. Whether that selectivity is optical, chemical or
inferential is not recoverable from the title, and no performance claim is asserted.
On the retry list.}
{10.1016/j.buildenv.2026.115281}

\paperentry{Raju et al.\ 2026}
{Investigations of black carbon aerosol sources and strengths over a rural-valley location
in Haryana, India}
{J Environ Sci Health A}
{Three years of seven-wavelength \textbf{AE33 aethalometry} (January 2020 -- December 2022)
at Panchgaon, a rural Aravali valley site at 285\,m AMSL in northwest India. Annual means:
eBC \textbf{4.06\,$\pm$\,2.15}, fossil-fuel eBCff \textbf{3.10\,$\pm$\,2.40}, biomass-burning
eBCbb \textbf{1.42\,$\pm$\,1.19}\,$\mu$g/m$^3$, with eBCff carrying \textbf{68\%} of
absorption at 880\,nm. Winter maxima (eBCff 4.19, eBCbb 2.51\,$\mu$g/m$^3$) and monsoon
minima (1.90, 0.52) track boundary-layer height and wet scavenging. The aethalometer-model
split rests on \textbf{assumed absorption \AA ngstr\"om exponents that are never reported},
and the eBCff/eBCbb ratio of \textbf{3.02\,$\pm$\,3.80} has a standard deviation larger
than its mean --- the apportionment is unstable even where the record is not. A three-year
reference-grade absorption series at a rural South Asian site is exactly the co-location
anchor low-cost networks there lack.}
{10.1080/10934529.2026.2734478}


\clearpage

\band{Exposure assessment \& modelling\hfill 9 records}{Amber}

\paperentry{Kim et al.\ 2026}
{Residential indoor air monitoring for long-term exposure assessment in susceptible
populations}
{Environ Health}
{130 older adults with idiopathic pulmonary fibrosis, Seoul Metropolitan Area, 2021--23.
Indoor hourly PM$_{2.5}$ averaged \textbf{14.90\,$\mu$g/m$^3$ (SD 36.12)} against
\textbf{20.44 (61.63)} outdoors, higher in winter (20.17) than summer (11.19) and by day
than night, with hourly patterns that contrast rather than track outdoors. Infiltration
factors from median indoor/outdoor ratios of sulfur, nickel and iron were
\textbf{51--70\%}. Source contrasts were large: pets 24.09 vs 13.35, humidifiers 19.19 vs
13.16, smokers 16.76 vs 11.79, cooking 15.14 vs 10.96\,$\mu$g/m$^3$. Every standard
deviation exceeds its mean, so the distribution is episodic and a seasonal mean is a poor
summary; five days per season is thin support for a long-term claim and no health estimate
is reported yet. The transferable piece is the non-contact protocol --- instruments
redeployed by postal mail, which is how an unattended network is actually serviced.}
{10.1186/s12940-026-01334-3}

\paperentry{Yu et al.\ 2026}
{Continental pollutants modulate organic nitrogen and light absorption of marine organic
aerosols}
{Atmos Chem Phys}
{Shipboard sampling over the Yellow and Bohai Seas across three seasons. Aerosol organic
nitrogen averaged \textbf{0.35\,$\pm$\,0.25\,$\mu$g\,N\,m$^{-3}$}, 4--60\% of total
nitrogen, highest in autumn. Anthropogenic secondary pollutants --- aged biomass burning
and secondary nitrate --- supplied \textbf{36--76\% of ON}, 46--83\% of the water-soluble
and 39--89\% of the water-insoluble fraction; in spring dust accounted for 55\% of ON and
above 80\% during a dust episode, while summer carried a marine biogenic component.
Organic-aerosol light absorption rose with both ON concentration \emph{and} ON/OC ratio,
so absorption capability per unit carbon increased with nitrogen enrichment rather than
tracking mass. Three cruises give no annual coverage and apportionment is receptor
modelling on a small campaign. The sensing consequence is direct: an optical instrument
calibrated on one absorption regime is misapplied when composition shifts.}
{10.5194/acp-26-13175-2026}

\paperentry{Sunyer-Cald\'u 2026}
{Beyond blood and urine: decoding the human chemical exposome across matrices and time}
{Environ Sci Technol}
{A Perspective arguing that high-resolution mass spectrometry can profile thousands of
chemical features per sample while study design has not kept pace --- no investigation has
combined external source monitoring, multi-matrix internal biomonitoring and toxicokinetic
normalisation in one longitudinal framework. Three elements are proposed: multi-matrix
sampling extending beyond blood and urine to accumulation archives, developmental matrices
and external media; a Unified Exposure Scale placing heterogeneous matrices on a common
toxicokinetically normalised dosimetric axis; and a two-track hazard-prioritisation
pipeline. No data, no validation, and the Unified Exposure Scale is circular where it
matters --- it needs toxicokinetic parameters that do not exist for the molecular dark
matter it is meant to rank. Logged for how air monitoring is being positioned inside the
Global Human Exposome Project.}
{10.1021/acs.est.6c06509}

\paperentry{Godri Pollitt 2026 \textnormal{---} metadata only}
{Beyond inhalation: air as an integrative matrix for exposome research}
{Environ Sci Technol}
{ACS deposited this DOI on 18 September without an abstract and no text was recoverable
from Crossref, Europe PMC or PubMed on per-DOI retry. Carried because it is the
\textbf{second air-as-exposome-matrix Perspective in this same issue}, alongside
Sunyer-Cald\'u, which makes the pair an editorial signal rather than a coincidence: air
sampling is being argued into the exposome programme as an integrative matrix rather than
merely an inhalation route. No argument or claim is asserted beyond the title.
On the retry list.}
{10.1021/acs.est.6c05073}

\paperentry{Sun et al.\ 2026}
{Nonlinear nitrate formation governed by aerosol water and VOC-sensitive oxidation in a
humid urban basin}
{Environ Sci Technol}
{Wintertime field observations in a humid, topographically confined Chinese urban basin,
combined with a box model and thermodynamic calculations. Nitrate became the dominant
PM$_{2.5}$ component during pollution episodes, co-varying with NO$_2$, O$_x$ and aerosol
liquid water content. Pathway-resolved modelling split formation almost evenly between
daytime OH+NO$_2$ (\textbf{52\%}) and nighttime heterogeneous N$_2$O$_5$ hydrolysis
(\textbf{46\%}), the latter rising sharply during episodes. Binning regressions located
thresholds at \textbf{RH 87\%} and \textbf{ALWC 76.5\,$\mu$g/m$^3$} above which formation
efficiency rose abruptly; VOCs exerted slightly stronger control than NO$_x$, with
oxygenated VOCs and alkenes dominant. One site, one winter, one box model --- the
thresholds are not validated out of sample. The RH threshold is directly load-bearing for
sensor correction: it is the same humidity range in which optical low-cost sensors diverge
from reference monitors.}
{10.1021/acs.est.6c08621}

\paperentry{Zhao et al.\ 2026}
{Unveiling the importance of organic hydroxy acids-related autoxidation pathway in isoprene
secondary organic aerosol formation}
{Environ Sci Technol}
{Year-long PM$_{2.5}$ measurements at \textbf{nine Pearl River Delta sites} quantifying six
isoprene-derived highly oxygenated molecules (iHOMs). Annual mean
\textbf{2.69\,$\pm$\,3.52\,ng/m$^3$}, higher in summer and autumn, dominated by
methyltartaric acids (\textbf{71\%}). Particle-phase tracer shares for the HOM, HMML and
IEPOX pathways were 10\%, 11\% and \textbf{79\%}; box-modelled gas-phase intermediates gave
25\%, 22\% and 53\%, so IEPOX dominates isoprene SOA in the polluted PRD with the other two
non-negligible in winter. Pathway partitioning tracked the ISOPO$_2$ isomer split. The
standard deviation exceeds the mean, so the annual average describes episodes rather than a
typical day, and tracer-to-SOA scaling relies on chamber yields measured well above ambient.
These are \textbf{nanogram-level organics} --- the composition layer no optical network sees.}
{10.1021/acs.est.6c06525}

\paperentry{Frueh et al.\ 2026}
{Source apportionment of volatile organic compounds in a Philadelphia fenceline community:
results from THRIVEair}
{J Expo Sci Environ Epidemiol}
{Year-long community-designed VOC campaign in South Philadelphia run with Philly Thrive,
collecting one-week integrated passive thermal-desorption samples of \textbf{38 VOCs at 20
sites}, July 2023 -- June 2024, around a former oil refinery under redevelopment.
Concentrations were higher in industrial than non-industrial and high- than low-traffic
areas. \textbf{Benzene, the community's declared concern, was low overall} (median
1.21\,$\mu$g/m$^3$, range 0.31--9.05), while (\textit{m,p})-xylenes emerged as the
industrial-area pollutant of concern; toluene:benzene and xylene:ethylbenzene ratios were
higher at industrial and high-traffic sites, indicating fresher emissions, and PMF
apportioned sources. \textbf{One-week integration resolves none of the short plumes
fenceline residents actually report.} The design is the template distributed PM networks
should copy: community-set siting and a declared pollutant of concern rather than
compliance checking.}
{10.1038/s41370-026-00976-2}

\paperentry{Konon et al.\ 2026 \textnormal{---} metadata only}
{Seasonal variation in volatile organic compound emissions and secondary organic aerosol
formation from unconventional oil and gas development}
{Atmos Environ}
{Same-day Elsevier deposit; nothing in Crossref, Europe PMC core, Europe PMC
\texttt{fullTextXML} or the OpenAlex per-work record. The title indicates a seasonally
resolved measurement campaign of VOC emissions from unconventional oil and gas development
and the secondary organic aerosol those emissions form. In scope as a precursor study:
\textbf{SOA is the fraction of basin PM$_{2.5}$ that no primary-particle emission inventory
captures}, and seasonality in the precursor flux propagates straight into any exposure
model built on annual-average emission factors. Carried forward for pickup.}
{10.1016/j.atmosenv.2026.122368}

\paperentry{Dou et al.\ 2026 \textnormal{---} metadata only}
{Volatile organic compounds in Chengdu: functional zone disparities and four-year dynamics
in a basin city of southwest China}
{Environ Pollut}
{Same-day Elsevier deposit; no abstract in Crossref, Europe PMC core, Europe PMC
\texttt{fullTextXML} or the OpenAlex per-work record. The title indicates a \textbf{four-year}
ambient VOC record for Chengdu resolved by urban functional zone --- the same Sichuan Basin
whose humidity-driven nitrate chemistry Sun et al.\ model in this issue. In scope as
precursor monitoring: functional-zone contrasts are the spatial scale distributed networks
are built to resolve, and a four-year baseline is long enough to separate trend from
meteorological variability. Carried forward for pickup.}
{10.1016/j.envpol.2026.129164}

\band{Burden, policy \& mitigation\hfill 1 record}{Coral}

\paperentry{Massimi et al.\ 2026}
{Urban green walls under traffic stress: elemental composition of PM deposition and plant
response}
{Environ Sci Technol}
{A nine-species green wall at a high-traffic urban site, followed for a year. PCA on the
insoluble fraction of Al, Cr, Cs, Fe, Mn, Mo, Rb, Sb, Sn, Ti and Zn identified
\textit{Photinia} $\times$ \textit{fraseri} as the highest and most stable retainer of
non-exhaust traffic particles and \textit{Pteris vittata} as the lowest, with SEM--EDS
showing \textbf{up to 100-fold higher retention} in the former. The methodological finding
travels furthest: \textbf{only insoluble elements traced deposition reliably}, because
soluble and whole-leaf fractions are dominated by endogenous pools of nutrients and of
tracers such as Cs and Rb, so conventional whole-leaf digestion overestimates deposition.
\textit{Photinia} showed marked seasonal physiological and transcriptional modulation;
\textit{Pteris} did not. One wall at one site with no paired ambient PM mass measurement,
so species are ranked against each other but no removal flux is quantified and no
air-quality benefit can be inferred.}
{10.1021/acs.est.6c10086}


\clearpage

% ==================================================================== TABLE
\section*{\sffamily\bfseries\color{Deep}Corpus register}
\vspace{-1.5mm}{\color{Amber}\rule{\linewidth}{1.1pt}}\vspace{2.5mm}

{\fontsize{6.7}{8.4}\selectfont
\setlength{\tabcolsep}{2.4pt}
\begin{longtable}{@{}>{\RaggedRight}p{25mm}>{\RaggedRight}p{41mm}>{\RaggedRight}p{27mm}%
>{\RaggedRight}p{23mm}>{\RaggedRight}p{18mm}>{\RaggedRight}p{15mm}@{}}
\toprule
\textsf{\textbf{First author}} & \textsf{\textbf{Focus}} & \textsf{\textbf{Journal}} &
\textsf{\textbf{Design}} & \textsf{\textbf{Metric}} & \textsf{\textbf{Region}}\\
\midrule
\endfirsthead
\toprule
\textsf{\textbf{First author}} & \textsf{\textbf{Focus}} & \textsf{\textbf{Journal}} &
\textsf{\textbf{Design}} & \textsf{\textbf{Metric}} & \textsf{\textbf{Region}}\\
\midrule
\endhead
\multicolumn{6}{@{}l@{}}{\cellcolor{Mist}\textsf{\textbf{\textcolor{Deep}{Neuro / mental health}}}}\\[0.6mm]
\href{https://doi.org/10.1001/jamanetworkopen.2026.25015}{Dales} & Wildfire-sourced PM$_{2.5}$ and ED visits for migraine and other primary headache & \textit{JAMA Netw Open} & Case-crossover & PM$_{2.5}$ (wildfire) & Canada\\
\href{https://doi.org/10.1093/ageing/afag215}{Glover} & Umbrella meta-review of urban environmental exposures and cognitive health & \textit{Age Ageing} & Umbrella review / meta-analysis & PM$_{2.5}$ & Global\\
\href{https://doi.org/10.1016/j.nsa.2026.107019}{Goossens} & In utero air pollution exposure and pre-/postnatal brain development & \textit{Neurosci Appl} & Narrative review & PM$_{2.5}$ / PM$_{10}$ / PM$_{\mathrm{coarse}}$ & Global\\
\href{https://doi.org/10.1016/j.isci.2026.116617}{Ma} & PM$_{2.5}$-induced neurotoxicity and depression: exposure risks and mechanisms & \textit{iScience} & Narrative review & PM$_{2.5}$ + components & Global\\
\href{https://doi.org/10.3760/cma.j.cn112150-20250804-00752}{Pan} & Maternal PM exposure and offspring neurodevelopment: effects and mechanisms & \textit{Zhonghua Yu Fang Yi Xue Za Zhi} & Narrative review & PM (bulk) & China\\
\addlinespace[0.9mm]
\multicolumn{6}{@{}l@{}}{\cellcolor{Mist}\textsf{\textbf{\textcolor{Teal}{Cardiovascular \& metabolic}}}}\\[0.6mm]
\href{https://doi.org/10.3760/cma.j.cn112338-20251120-00843}{Huang} & PM$_{2.5}$ and dyslipidaemia in rural Chinese children; dietary effect modification & \textit{Zhonghua Liu Xing Bing Xue Za Zhi} & Cross-sectional (multistage) & PM$_{2.5}$ & China\\
\href{https://doi.org/10.32394/pe/218966}{Kazmierska} & Air pollution and cardiovascular disease: review of current data & \textit{Przegl Epidemiol} & Narrative review & PM$_{2.5}$ / PM$_{10}$ & Global\\
\href{https://doi.org/10.1007/s00125-026-06814-2}{Xi} & Air pollution, built environment and temperature vs. prediabetes phenotypes & \textit{Diabetologia} & Prospective cohort & PM$_{2.5}$ & Germany\\
\href{https://doi.org/10.1161/JAHA.125.046917}{Yang} & Short-term PM$_{2.5}$ and hospitalisation/ED use in atrial fibrillation & \textit{J Am Heart Assoc} & Case-crossover & PM$_{2.5}$ (wildfire-influenced) & USA\\
\addlinespace[0.9mm]
\multicolumn{6}{@{}l@{}}{\cellcolor{Mist}\textsf{\textbf{\textcolor{Sky}{Reproductive \& developmental}}}}\\[0.6mm]
\href{https://doi.org/10.1016/j.jes.2025.10.043}{Chen} & Prenatal air pollution, fetal growth trajectory phenotypes and child neurodevelopment & \textit{J Environ Sci (China)} & Prospective cohort & PM$_{2.5}$ / PM$_{10}$ & China\\
\href{https://doi.org/10.1016/j.ecoenv.2026.120566}{Huang} & Joint prenatal bisphenol + air-pollutant mixture exposure and preschool IQ & \textit{Ecotoxicol Environ Saf} & Prospective cohort (mixtures) & PM$_{2.5}$ / PM$_{10}$ & China\\
\href{https://doi.org/10.1007/s40572-026-00552-8}{Peterson} & Ambient air pollution, PAHs and pubertal development: critical review & \textit{Curr Environ Health Rep} & Critical review & PM$_{2.5}$ components & Global\\
\href{https://doi.org/10.1016/j.envres.2026.125313}{Xie} & Long-term PM$_{2.5}$ and timing of pubertal development in boys and girls & \textit{Environ Res} & Prospective cohort & PM$_{2.5}$ & China\\
\addlinespace[0.9mm]
\multicolumn{6}{@{}l@{}}{\cellcolor{Mist}\textsf{\textbf{\textcolor{Sage}{Respiratory \& allergic}}}}\\[0.6mm]
\href{https://doi.org/10.1007/s00484-026-03289-z}{Bayar Muluk} & Climate change and otolaryngology: airway disease impacts and emission reduction & \textit{Int J Biometeorol} & Narrative review & PM$_{2.5}$ & Global\\
\href{https://doi.org/10.4103/mgr.MEDGASRES-D-25-00164}{Liao} & Short-term air pollution, respiratory symptoms and lung function in asthma & \textit{Med Gas Res} & Panel study & PM$_{2.5}$ / PM$_{10}$ & China\\
\href{https://doi.org/10.3390/epidemiologia7040097}{Sanchez de Prada} & Ambient PM and asthma-coded admissions in children 0-3 y, Spain & \textit{Epidemiologia} & Ecological panel & PM$_{2.5}$ / PM$_{10}$ & Spain\\
\href{https://doi.org/10.1016/j.jhazmat.2026.143076}{Yang} & Allergenic vs. toxic risk of airborne seed hairs compared with ambient PM$_{2.5}$ & \textit{J Hazard Mater} & Field measurement & PM$_{2.5}$ + bioaerosol & China\\
\addlinespace[0.9mm]
\multicolumn{6}{@{}l@{}}{\cellcolor{Mist}\textsf{\textbf{\textcolor{Violet}{Mechanistic toxicology}}}}\\[0.6mm]
\href{https://doi.org/10.1016/j.jes.2025.12.074}{Ding} & PM$_{2.5}$ promotes vascular calcification via miR-27b-5p/HES1 signalling & \textit{J Environ Sci (China)} & In vitro + organoid + murine & PM$_{2.5}$ & China\\
\href{https://doi.org/10.1007/s00360-026-01698-5}{Maraschi} & Osmoregulatory disruption and metal bioaccumulation from settleable atmospheric PM & \textit{J Comp Physiol B} & Controlled exposure (ecotox) & Settleable APM & Brazil\\
\href{https://doi.org/10.1016/j.jes.2025.11.018}{Wu} & Macrophage NLRP3 inflammasome in real-ambient PM-induced hepatic glucose dysregulation & \textit{J Environ Sci (China)} & Real-ambient murine + co-culture & Real-ambient PM & China\\
\addlinespace[0.9mm]
\multicolumn{6}{@{}l@{}}{\cellcolor{Mist}\textsf{\textbf{\textcolor{Clay}{Occupational \& indoor}}}}\\[0.6mm]
\href{https://doi.org/10.3389/fmedt.2026.1796828}{Chen} & Nested suction instrument to cut intraoperative surgical-smoke PM$_{2.5}$ & \textit{Front Med Technol} & Bench / device experiment & PM$_{2.5}$ (surgical smoke) & China\\
\href{https://doi.org/10.1093/annweh/wxag055}{Cole} & Determinants of respirable crystalline silica exposure in Australian tunnelling & \textit{Ann Work Expo Health} & Observational exposure study & Respirable dust / RCS & Australia\\
\href{https://doi.org/10.1016/j.jhazmat.2026.143081}{Ding} & Stage-resolved hair Zn-isotope dynamics as a firefighting combustion-exposure tracer & \textit{J Hazard Mater} & Repeated-measures biomarker & Combustion-derived PM & China\\
\href{https://doi.org/10.3390/jox16040119}{Ramsden} & Xenobiotic hazards in aircraft cabin air & \textit{J Xenobiot} & Review + risk economics & Ultrafine particles & Global\\
\addlinespace[0.9mm]
\multicolumn{6}{@{}l@{}}{\cellcolor{Mist}\textsf{\textbf{\textcolor{Amber}{Sensing, forecasting \& instrumentation}}}}\\[0.6mm]
\href{https://doi.org/10.1016/j.jenvman.2026.130515}{Shi} & Adaptive spatiotemporal graph transformer for multi-site multi-step PM$_{2.5}$ forecasting & \textit{J Environ Manage} & Deep-learning model & PM$_{2.5}$ & China / India\\
\href{https://doi.org/10.3390/bios16070380}{Tsuchiya} & Seasonal and indoor-outdoor characteristics of fluorescent aerosol particles in offices & \textit{Biosensors} & Real-time sensor campaign & Bioaerosol / coarse & Japan\\
\addlinespace[0.9mm]
\multicolumn{6}{@{}l@{}}{\cellcolor{Mist}\textsf{\textbf{\textcolor{Amber}{Exposure assessment \& modelling}}}}\\[0.6mm]
\href{https://doi.org/10.1029/2025GH001621}{Solorzano-Araque} & Multiscale exposure modelling: LOTOS-EUROS + mobility surveys + particle toxicity & \textit{GeoHealth} & CTM + mobility integration & PM$_{2.5}$ / PM$_{10}$ & Colombia\\
\addlinespace[0.9mm]
\multicolumn{6}{@{}l@{}}{\cellcolor{Mist}\textsf{\textbf{\textcolor{Coral}{Burden, policy \& mitigation}}}}\\[0.6mm]
\href{https://doi.org/10.1002/hsr2.72886}{Karamzad} & Ischaemic heart disease burden in MENA 1990-2021; PM attributable fraction 28.2\% & \textit{Health Sci Rep} & GBD secondary analysis & Ambient PM & MENA\\
\href{https://doi.org/10.1371/journal.pone.0347989}{Rees} & HEPA filtration to reduce respiratory infection in care homes & \textit{PLoS One} & RCT process evaluation & Filtered PM & UK\\
\href{https://doi.org/10.1016/j.jes.2026.01.009}{Shan} & Policy-driven emission reductions and post-pandemic rebound, high-resolution inventories & \textit{J Environ Sci (China)} & Emission inventory & PM emissions & China\\
\addlinespace[0.9mm]
\multicolumn{6}{@{}l@{}}{\cellcolor{Mist}\textsf{\textbf{\textcolor{Slate}{Other clinical endpoints}}}}\\[0.6mm]
\href{https://doi.org/10.3390/dj14070440}{Lupita} & PM$_{2.5}$, preventive-care capacity and caries/edentulism burden across the EU & \textit{Dent J} & Ecological panel & PM$_{2.5}$ & EU-27\\
\href{https://doi.org/10.1016/j.jdin.2026.06.008}{Park} & Long-term PM exposure and risk of androgenetic alopecia & \textit{JAAD Int} & Nationwide cohort & PM$_{2.5}$ / PM$_{10}$ & South Korea\\
\href{https://doi.org/10.3171/2026.3.PEDS2615}{Sun} & Environmental exposure and area deprivation vs. paediatric hydrocephalus outcomes & \textit{J Neurosurg Pediatr} & Registry cohort + geospatial & PM$_{2.5}$ (satellite-derived) & USA\\
\addlinespace[0.9mm]
\bottomrule
\end{longtable}}


\vspace{4mm}
\begin{keybox}
{\sffamily\bfseries\footnotesize\color{Deep}Provenance and method}\\[1.2mm]
{\fontsize{7.2}{8.8}\selectfont
Entry window \textbf{18 September 2026}, single day, continuous with the previous issue.
This issue was built in \textbf{two passes}: a 10:33 CDT pass that saw roughly ten hours of
the window and admitted 17 records, and a 23:0x re-harvest of the same window that admitted
13 more. The second pass is what makes the window complete; see the run note above.
Legs run in the second pass: \textbf{PubMed} E-utilities on the health and instrumentation
axes as separate queries (\textbf{12} health, 0 sensing); \textbf{Europe PMC}
\texttt{CREATION\_DATE} (\textbf{46}); \textbf{Crossref by ISSN} (\textbf{53}, again the
highest-yield leg). \textbf{OpenAlex} returned HTTP 429 and \textbf{arXiv} HTTP 406, both
known-broken legs. The \textbf{PubMed connector} added \textbf{no} net on-topic record: of
its twelve PMIDs, three were unseen on the PMID key, but two of those three are records this
issue already carries --- the morning pass admitted them from Europe PMC with no PMID
attached, so only the DOI key caught them. The one genuinely new record was an x-ray
air-kerma Monte Carlo study, rejected as off-topic.
\textbf{Nine} Europe PMC records were dropped at the source as retro-index artefacts
(AGRICOLA and PMC records returned inside the window whose first publication dates precede
the issue month --- the gate hand-applied on 18 September now lives in the screen script).
The remaining 102 raw records deduplicated to 97, of which 25 were already carried, leaving
72 fresh candidates of which \textbf{13 were admitted}. Across both passes, \textbf{132 fresh candidates were screened and 30 admitted}.
\textbf{Consensus} returned 10 sensor-calibration papers and admitted \textbf{none} for the
third consecutive issue: 5 matched archive titles and all 5 survivors failed a Crossref
\texttt{created} check (2026-01-02 to 2026-06-15), including a 1,013-sensor global PurpleAir
calibration whose deposit date is 15 June. A \textbf{new abstract-recovery route was tested
and failed}: the free OpenAlex per-work endpoint returns an abstract only where one was
deposited to Crossref, so it recovers nothing from the publishers that withhold them.
\textbf{Eight records are carried on metadata alone} and assert no finding. All 30 DOIs
were resolved against Crossref before build (0 fail, 0 warn).
Subtopic, design, particle-metric and geography labels are assigned from title, abstract,
keywords and MeSH; assignment is single-label by dominant contribution and is a judgement,
not a controlled vocabulary. Effect estimates are transcribed verbatim and use
heterogeneous exposure contrasts --- see the forest-plot caption.
Abstracts remain the copyright of their respective publishers.\par}
\end{keybox}

\end{document}
